% ============================================================
%  组会汇报 Beamer 模板 —— 基于组会PPT.tex（Madrid + 二月兰紫）
%  编译方式: XeLaTeX，需编译两次
% ============================================================
\documentclass[aspectratio=169]{beamer}

\mode<presentation> {
	\usetheme{Madrid}
	\setbeamertemplate{navigation symbols}{}
	\setbeamertemplate{footline}[frame number]
}

% ---------- 二月兰紫配色 ----------
\definecolor{orchidpurple}{RGB}{146,39,143}
\definecolor{orchiddark}{RGB}{115,30,112}
\definecolor{orchidlight}{RGB}{230,210,230}

\setbeamercolor{structure}{fg=orchidpurple}
\setbeamercolor{palette primary}{fg=white,bg=orchidpurple}
\setbeamercolor{palette secondary}{fg=white,bg=orchiddark}
\setbeamercolor{palette tertiary}{fg=white,bg=orchidpurple}
\setbeamercolor{palette quaternary}{fg=white,bg=orchiddark}
\setbeamercolor{titlelike}{fg=white,bg=orchidpurple}
\setbeamercolor{frametitle}{fg=white,bg=orchidpurple}
\setbeamercolor{block title}{fg=white,bg=orchidpurple}
\setbeamercolor{block body}{fg=black,bg=orchidlight!30}
\setbeamercolor{item}{fg=orchidpurple}
\setbeamercolor{section in toc}{fg=orchiddark}

\usepackage{graphicx}
\usepackage{booktabs}
\usepackage[UTF8,noindent]{ctexcap}
\usepackage[bookmarks=true]{hyperref}
\usepackage{amsmath,amssymb,bm}

% ---------- 快捷命令 ----------
\newcommand{\T}{^{\mathsf{T}}}
\newcommand{\R}{\mathbb{R}}
\newcommand{\mat}[1]{\bm{#1}}
\newcommand{\tbf}[1]{\textbf{#1}}

% ============================================================
%	以下为可修改区：标题 / 作者 / 单位 / 日期
% ============================================================

\title[汇报主题（短标题）]{汇报主题（完整标题）}
\author{汇报人}
\institute[南京理工大学]
{
	南京理工大学 \\
	\medskip
	\textit{邮箱@njust.edu.cn}
}
\date{YYYY年M月}

% ============================================================
%	以下为正文示例，请按需替换
% ============================================================

\begin{document}

% ---------- 标题页 ----------
\begin{frame}
	\titlepage
\end{frame}

% ---------- 目录 ----------
\section*{目录}
\begin{frame}
	\frametitle{汇报提纲}
	\tableofcontents
\end{frame}

%-----------------------------------
%	一、请替换为你的章节名
%-----------------------------------
\section{第一章：章节标题}

\begin{frame}
	\frametitle{双栏：要点 + 流程图}
	\begin{columns}[T]
		\column{0.55\textwidth}
		\begin{itemize}
			\item \tbf{要点一}：说明内容
			\item \tbf{要点二}：说明内容
			\item \tbf{要点三}：说明内容
		\end{itemize}

		\column{0.45\textwidth}
		\vspace{0.3cm}
		\begin{block}{流程块}
			\begin{enumerate}
				\item 第一步
				\item 第二步
				\item 第三步
			\end{enumerate}
		\end{block}
	\end{columns}
\end{frame}

\begin{frame}
	\frametitle{公式展示}
	\begin{itemize}
		\item 转置符号：$\mat{x}\T$
		\item 实数集：$\R^n$
		\item 动力学方程：$\mat{M}(\mat{q})\ddot{\mat{q}} + \mat{C}(\mat{q},\dot{\mat{q}})\dot{\mat{q}} + \mat{g}(\mat{q}) = \bm{\tau}$
	\end{itemize}
\end{frame}

%-----------------------------------
%	二、请替换为你的章节名
%-----------------------------------
\section{第二章：章节标题}

\begin{frame}
	\frametitle{图文混排}
	\begin{columns}[T]
		\column{0.5\textwidth}
		\tbf{技术亮点}
		\begin{enumerate}
			\item \textbf{亮点一}：说明
			\item \textbf{亮点二}：说明
			\item \textbf{亮点三}：说明
		\end{enumerate}

		\column{0.5\textwidth}
		\tbf{主要收获}
		\begin{itemize}
			\item 收获一
			\item 收获二
			\item 收获三
		\end{itemize}
	\end{columns}
\end{frame}

\begin{frame}
	\frametitle{插入图片}
	\begin{center}
		\includegraphics[width=0.6\textwidth]{example-image}\\
		\small 图注：请替换为你的图片
	\end{center}
\end{frame}

\begin{frame}
	\frametitle{插入表格}
	\begin{table}
		\centering
		\begin{tabular}{lcc}
			\toprule
			列名 & 数据A & 数据B \\
			\midrule
			行1 & 数值 & 数值 \\
			行2 & 数值 & 数值 \\
			\bottomrule
		\end{tabular}
	\end{table}
\end{frame}

% ---------- 结束页 ----------
\begin{frame}
	\centering
	\Huge{感谢聆听！}\\[1cm]
	\Large{敬请各位老师、同学批评指正}
\end{frame}

\end{document}
